\input{ieee-preamble}
\hypersetup{
  pdftitle={Entropy-Minimizing Control of a Simulated Helicopter},
  pdfauthor={Santiago Restrepo},
  pdfsubject={Native C++ nonlinear predictive control with explicit thermodynamic accounting},
  pdfkeywords={helicopter, physical entropy generation, nonlinear MPC, simulation}
}
\title{Entropy-Minimizing Control of a Helicopter}
\author{\IEEEauthorblockN{Santiago Restrepo}
\IEEEauthorblockA{\href{mailto:contact@waajacu.com}{contact@waajacu.com}\quad
\url{https://waajacu.com/}\\[3pt]
\textit{An explicitly parameterized simulation study}}}
\date{}
\begin{document}
\maketitle
\begin{abstract}
Native C++ predictive control minimizes physical entropy generation in a
simulated electric helicopter with rotor and thermal dynamics. Airborne
optimization, PID fallback, and PID landing support are reported separately.
Parameters are assumed; a matched PID baseline supports evaluation.
\end{abstract}
\begin{IEEEkeywords}
Helicopter, entropy, NMPC.
\end{IEEEkeywords}
\balance
\section{Fix the Vehicle and Mission}
The 2\,kg aircraft has main/tail rotor radii 0.55/0.11\,m and governed
speeds 180/650\,rad/s. Its 23 states contain
position, velocity, quaternion, body rates, actual pitches, induced velocities,
flapping angles, and motor temperatures. Inputs are physical pitch commands:
\begin{equation}
\dot x=f(x,u,w),\qquad
u=[\theta_0,\theta_{1c},\theta_{1s},\theta_t]^{\mathsf T}.
\end{equation}
Both controllers share a 70\,s takeoff--waypoint--landing mission
and exact simulated state/wind. The simulator supplies a spooled trim for each
actual mass; controllers do not discover that trim. MuJoCo advances rigid-body/contact dynamics
at 0.01\,s; rotor and thermal states use RK4. Both controllers share a causal
force/velocity mass estimate, assuming inertia scales with mass. The initial
0.1\,s trim observation hold is included in mission accounting.

\section{Close the Physical Balances}
The boundary includes motor heating, drivetrain loss, aerodynamic
dissipation, and wake relaxation at ambient $T_0=293.15$\,K. For rotor $r$,
near-hover blade-element and inflow relations give~\cite{johnson2015}
\begin{equation}
P_{a,r}=F_r V_{a,r}+\kappa_r F_r v_{i,r}
       +P_{\mathrm{profile},r}+\dot E_{i,r}.
\end{equation}
Near-wake storage is $E_{i,r}=M_{i,r}v_{i,r}^2/2$.
Shaft power includes body rotation. Copper, iron, drivetrain, and braking losses
contribute to motor heat $L_r$. With heat rejection $H_r=h_r(T_r-T_0)$,
\begin{equation}
\begin{aligned}
C_r\dot T_r&=L_r-H_r,\\
\sigma&=\sum_r\left[\frac{L_r}{T_r}
 +H_r\left(\frac1{T_0}-\frac1{T_r}\right)\right]
 +\frac{D_a}{T_0}.
\end{aligned}
\end{equation}
$D_a$ comprises modeled wake, profile, translational-drag and angular-damping
dissipation. Wind and contact work are separate mechanical ports. Electrical
input, storage, heat rejection, and entropy are independently audited.

\section{Optimize the Physical Objective}
Native CasADi/IPOPT~\cite{casadi} seeks local solutions by multiple shooting with 20
intervals of 0.1\,s and a 2\,s prediction horizon:
\begin{equation}
\begin{aligned}
\min_U\ J&=\int_0^{2\,\mathrm{s}}\!\sigma(x,u,w)\,dt+\Phi(x_N),\\
\Phi&=\frac1{T_0}\sum_r\left[E_{i,r}
+C_r\left(T_r-T_0-T_0\ln\frac{T_r}{T_0}\right)\right].
\end{aligned}
\end{equation}
$\Phi$ accounts for wake relaxation and passive motor cooling.
Tracking is constrained, with no tracking-error cost. Bounds cover commands,
flight, electrical and thermal states, and terminal settling; they do not
prove stability or recursive feasibility.
Independently reintegrated feasible solutions drive airborne flight; rejected
solves use PID fallback. Before skid contact, geometry triggers explicit PID
landing support. This is not contact-optimal control. Playback may be
slower than wall time; optimization runs on CPU.

\section{Evidence, Limits, and Project}
Five matched 70\,s wind/mass cases and a disabled-control test pass, with no
audited physical-limit violations. NMPC supplies 90.57--93.29\% of all mission
decisions; trim observation, fallback, and PID landing remain counted.
Committed entropy is 0.12--0.20\% lower than PID, exceeding audited accounting
and endpoint uncertainty. NMPC tracks less closely within the shared corridor.
Four handovers miss the tighter settling target but pass final landing checks.
Balance, cooldown, identification, numerical and rejection tests also pass.

The model omits servo electrical losses, rotor-speed transients, stall,
vortex-ring flow, and detailed ground aerodynamics. All solves miss the
0.1\,s deadline; no hardware efficiency is established. Exergy destruction rate is $T_0\sigma$~\cite{philipp2024};
matched endpoints can equate entropy and dissipated-energy savings.

\begingroup\raggedright
Code, assumed parameters, setup, viewer, and evaluation evidence:\\
\url{https://github.com/savethebeesandseeds/Universalis-Materiae-Dexteritas}\\
Directory: \texttt{projects/helicopter}.\par
\endgroup
\begingroup\raggedright
\bibliographystyle{IEEEtran}
\bibliography{references}
\endgroup
\end{document}
